% AUTORIDAD DE FUENTE: Exposé IX, §§8.2--8.3, líneas 2159--2188;
% índices de escaneo 427--428, folios 416--417. La única diferencia
% material D/C se conserva reversiblemente mediante \sisi. La prima de
% A'_K en la polarización se ha restituido como corrección transcriptiva
% inequívoca confirmada por la impresión.

\medskip
\noindent
8.2. Las hipótesis son las de 8.1. Es claro que la extensión (8.1.1)
depende funtorialmente del esquema abeliano $A_K$, y que su formación
conmuta con todo morfismo de rasgos $\widetilde S\longrightarrow S$.
Elijamos una polarización de $A_K$, y por tanto una isogenia
$A_K\longrightarrow\sisi{A_{K'}}{A'_K}$. Esta induce un homomorfismo de (8.1.1) a
(8.1.6), que se ve inmediatamente, considerando los $T_\ell$, que es
una isogenia término a término. Del mismo modo, utilizando 7.5 y el
descenso, se ve que toda biextensión por $\mathbb{G}_{mK}$ de esquemas
abelianos $A_K$ y $\overline A_K$ sobre $K$ define una biextensión de
$(B_K,\overline B_K)$ por $\mathbb{G}_{mK}$, donde $B_K$ y
$\overline B_K$ denotan los cocientes abelianos de las extensiones de
Raynaud correspondientes a $A_K$ y $\overline A_K$, respectivamente.
Cuando $W$ es una dualidad perfecta, también lo es la biextensión
correspondiente de $(B_K,\overline B_K)$. En particular, con la notación
de 7.6, los esquemas abelianos $B_K$ y $B'_K$ son canónicamente duales
entre sí.

La correspondencia entre las biextensiones sobre $(A_K,\overline A_K)$
y las biextensiones sobre $(B_K,\overline B_K)$ tiene propiedades de
functorialidad evidentes al variar $A_K$ y $\overline A_K$. Estas se
deducen por descenso de las propiedades análogas enunciadas en 7.5, y
dejamos su formulación al lector. También dejamos al lector la
compatibilidad de la construcción precedente con todo cambio de rasgos
$\widetilde S\longrightarrow S$.

\medskip
\noindent
8.3. Volvamos al grupo $A_{K,K'}^{\natural}$ y estudiemos su variación
con $K'$. Es claro que dos extensiones isomorfas $K'$ y $K'_1$ de
$\sisi{K'}{K}$ dan grupos canónicamente isomorfos, por lo que basta
considerar las subextensiones $K'$ de la clausura separable
$\overline K$ de $K$. Si $K'\hookrightarrow K''$ son dos de tales
subextensiones de Galois finitas de $\overline K$, se deduce
inmediatamente de 3.3 que existe una \underline{inmersión abierta}

%%%%%%%%%%  scan idx 428  |  folio 417  %%%%%%%%%%

\begin{equation}\tag{8.3.1}
A^{\natural}_{K,K'} \hookrightarrow A^{\natural}_{K,K''}
 \qquad ,
\end{equation}
que identifica el primer grupo con un subgrupo abierto del segundo. Al
variar $K'$, los $A^{\natural}_{K,K'}$ forman así un sistema inductivo
de grupos lisos de tipo finito sobre $K$, cuyos morfismos de transición
son inmersiones abiertas. Si ponemos

\begin{equation}\tag{8.3.2}
\Phi_{K,K'}=A^{\natural}_{K,K'}
 /{A^{\natural}_{K,K'}}^{o} \qquad ,
\end{equation}
los $\Phi_{K,K'}$ forman a su vez un sistema inductivo de grupos finitos
étale sobre $K$, que se estudiará con mayor detalle más adelante.

Pondremos

\begin{equation}\tag{8.3.3}
A^{\flat}_{K}=\varinjlim_{K'} A^{\natural}_{K,K'} \qquad ,
\end{equation}
donde el límite se toma en la categoría de los haces fppf sobre $K$.
Como los morfismos de transición de este sistema inductivo son
inmersiones abiertas, es claro además que este límite está representado
por un esquema en grupos sobre $K$ (en general no cuasicompacto), que
denotaremos naturalmente por el mismo símbolo $A^{\flat}_{K}$. Para
este esquema en grupos se tiene

\begin{equation}\tag{8.3.4}
{A_{K}^{\flat}}^{o}=A_{K}^{\natural o} \quad , \qquad
A^{\flat}_{K}/{A_{K}^{\flat}}^{o}
 \simeq \varinjlim_{K'}\Phi_{K,K'} \qquad .
\end{equation}
Este último grupo se determinará en 11.10: es canónicamente isomorfo a
$\underline{M}_{K}\otimes\mathbb{Q}/\mathbb{Z}$, con la notación de
(8.1.8).
